% Applied Software Development Lab -- Lab Handout 10: Now with Moving Pictures!
% Copyright (c) Michael Lehn. Licensed under CC BY-SA 4.0.
%
% Region coordinates (bp from the top of the page) follow the page layout of
% the original German handout "Laborblatt 10".
\documentclass{aslab}
\sethandout{10}{Now with Moving Pictures!}
\graphicspath{{../figures/handout10/}{../figures/shared/}}

% syntax colors of the code listings (used via escapechar=`)
\definecolor{codemaroon}{HTML}{970E53}
\newcommand*\Kw[1]{\textcolor{codeorange}{#1}}  % keywords
\newcommand*\Cy[1]{\textcolor{codecyan}{#1}}    % constants, names
\newcommand*\Nu[1]{\textcolor{codered}{#1}}     % numbers, strings
\newcommand*\Mg[1]{\textcolor{codemagenta}{#1}} % import
\newcommand*\Cm[1]{\textcolor{codepurple}{#1}}  % comments
\newcommand*\Bl[1]{\textcolor{aslabblue}{#1}}   % comments (p. 5), misc.
\newcommand*\St[1]{\textcolor{codemaroon}{#1}}  % string (p. 3)
% code lines keep the region pitch even with tall or deep glyphs
\makeatletter
\lst@AddToHook{Init}{\lineskiplimit=-\maxdimen\relax}
\makeatother
% "•" lists with the text indented by 9bp
\newlist{dotlist}{itemize}{2}
\setlist[dotlist]{label={\textbullet},leftmargin=9bp,labelsep=*,align=left,
  labelindent=0pt}
% numbered tasks with a regular (non-bold) number
\newlist{plaintasks}{enumerate}{1}
\setlist[plaintasks]{label={\arabic*.},leftmargin=18bp,labelsep=*,
  align=left,labelindent=0pt}
% (a), (b) lists with a bold label
\newlist{boldsteps}{enumerate}{1}
\setlist[boldsteps]{label={\bfseries(\alph*)},leftmargin=18bp,labelsep=*,
  align=left,labelindent=0pt}

\begin{document}
%% ===================================================================== 1 ==
\maketitleblock

\begin{region}{124.0}
\heading{Learning objectives}
\end{region}

\begin{region}[pitch=13]{138.5}
{\bfseries\color{aslabblue}
\begin{dashlist}
\item Refactoring (cleaning up is always part of it)
\item Catching exceptions
\item Avoiding data congestion when reading the serial interface
\item And above all: bitmaps and new game features
\end{dashlist}\par}
\end{region}

\begin{region}{204.5}
\heading{Today's goal}
\end{region}

\begin{region}[pitch=13.2]{219.5}
\begin{dashlist}
\item The controller should be connectable and usable even while the game
  is already running. But it should also be possible to remove it while the
  game is running. This is already supported in part -- but only in part.
\item If you turn the potentiometer quickly and use the push button at the
  same time, the game may ``lag'' (yes, that's really what people say!). The
  data you receive from the controller piles up. Even if you suddenly stop
  turning the potentiometer, the paddle keeps moving for a while.
\item Of course, a game also needs sound, e.g., when the ball bounces off,
  when a spaceship is shot down, etc.
\item A rocket should look like a rocket, a spaceship like a spaceship. In
  short: some things should be displayed with bitmaps and not just with
  geometric objects such as circles or rectangles.
\item And then, of course, you need lots of new features, for example
  spaceships that can drop bombs, bricks that have to be hit several times,
  etc. And maybe you shouldn't have an unlimited supply of ammunition
  either.
\item This handout uses a few examples to show how such new features can be
  implemented -- so that afterwards you can also implement your own ideas
  on your own.
\end{dashlist}
\end{region}

\begin{region}[pitch=13.18]{446.1}
We will tackle these things in a seemingly arbitrary order. This is meant to
show you what a typical process looks like in practice: at first, you only
intend to add one improvement -- and then one thing leads to another. You
discover places in the code that should be reorganized so that new features
fit in better. You find bugs that were hidden so far. And you get ideas for
the next feature that you didn't even have on your radar.
\end{region}

\placegraphic{148.9}{555.3}{297.5}{game-screenshot.png}

\nextpage
%% ===================================================================== 2 ==
\begin{region}{69.4}
\heading{Preparation: part 1}
\end{region}

\begin{region}[pitch=13.3]{88.1}
As usual, we work with two terminals: \textbf{left}: editing source code,
\textbf{right}: compiling code, testing the game. Clone the repository with
\end{region}

\begin{region}{119.2}
\code{git clone https://gitlab.uni-ulm.de/USERNAME/asp-ulm}
\end{region}

\begin{region}{136.9}
Then, in both terminals, change into the directory with the Python code:
\end{region}

\begin{region}{154.7}
\code{cd asp-ulm/python}
\end{region}

\placegraphic{280.9}{170.6}{257.6}{breadboard-setup.png}

\begin{region}[width=216]{177.4}
\heading{Preparation: part 2}
\end{region}

\begin{region}[width=216,pitch=13.17]{195.1}
You should build the familiar circuit with push button and potentiometer on
the breadboard and connect it correctly to the microcontroller. The push
button is connected to digital input pin~2, the potentiometer to analog
input A0. On the right, you can see a picture of the complete setup.
\end{region}

\begin{region}{305.8}
\heading{Preparation: part 3}
\end{region}

\begin{region}[pitch=13.45]{323.5}
In the first step today, the rockets should no longer be drawn as simple
filled circles but by an image. So that you don't have to draw the images
yourselves, I have created a Git repository with additional material.
\textbf{However, it is a bit tricky} to add these files to your own
repository in the end. Here are the instructions:
\end{region}

\begin{region}[pitch=13.3]{403.9}
\begin{tasks}
\item Make sure that you are in the subdirectory of your project. In the
  terminal, enter \code{pwd}. This prints the current directory.\\
  On the lab computers, this should be \code{/home/asp/asp-ulm/python}.
\end{tasks}
\end{region}

\begin{region}[pitch=17.8]{443.8}
\begin{tasks}[start=2]
\item Clone my repository with the additional material:\\
  \code{git clone https://github.com/michael-lehn/asp-ulm-data.git}
\end{tasks}
\end{region}

\begin{region}[x=74.7,right=0,pitch=14.2]{479.5}
This creates a new subdirectory \code{asp-ulm-data}.\\
\textbf{Confirm this} by displaying the content of the current directory with
\code{ls}.
\end{region}

\begin{region}[pitch=13.5]{507.3}
\begin{tasks}[start=3]
\item Rename the folder to \code{data}: \code{mv asp-ulm-data data}\\
  Use \code{ls} again to display what is in the current directory, and
  confirm that the subfolder is now called \code{data}.
\end{tasks}
\end{region}

\begin{region}[pitch=17.8]{547.6}
\begin{tasks}[start=4]
\item Now comes the delicate part, because something gets deleted. First
  use\\
  \code{ls -a data}
\end{tasks}
\end{region}

\begin{region}[x=74.7,right=0,pitch=13.25]{583.0}
to display everything that is in the folder. It should contain useful things
such as image and sound files, but also a directory \code{.git}. And exactly
this \code{.git} has to go! It contains administrative files because this
directory was originally cloned from my repository. You don't want to drag
this stuff along.\\
To do so, run:
\end{region}

\begin{region}[x=74.7,right=0,pitch=17.7]{653.9}
\code{rm -rf data/.git}\\
If you then run\\
\code{ls -a data}
\end{region}

\begin{region}[x=74.7,right=0]{707.2}
again, \code{.git} should be gone.
\end{region}

\begin{region}[pitch=17.8]{720.5}
\begin{tasks}[start=5]
\item Now add the folder with its entire content to your own repository:\\
  \code{git add data/*}
\end{tasks}
\end{region}

\begin{region}[x=74.7,right=0,pitch=13.5]{751.8}
\code{git commit -a -m "images and sound files"}\\
\code{git push}
\end{region}

\nextpage
%% ===================================================================== 3 ==
\begin{region}{69.4}
\heading{Step 1: Finally, here we go -- with images!}
\end{region}

\begin{region}[pitch=13.3]{87.4}
Connect the microcontroller and run \code{usb_check}.\\
Then open the file \code{ship.py} with the class \code{Ship}, which is
responsible for the spaceships.\\
They should no longer be drawn as simple filled circles, but look like in
Space Invaders.
\end{region}

\begin{region}{145.2}
In the constructor (i.e., the \code{__init__} method), add the following new
line:
\end{region}

\begin{region}{163.1}
\begin{codeblock}[escapechar=`]
self.img = pygame.image.load(`\St{"data/space\_invader0.png"}`)
\end{codeblock}
\end{region}

\begin{region}[pitch=13.35]{180.7}
This gives every Ship object a new attribute img, which is initialized with
an image object. And the image object was created by loading the file
\code{data/space_invader0.png}.
\end{region}

\begin{region}[pitch=13.3]{225.2}
Now this image should be drawn instead of the circle. So change the method
\code{draw} so that it contains only this line:
\end{region}

\begin{region}{256.3}
\tcode{pygame.Surface.blit(Game.screen, self.img, (self.x, self.y))}
\end{region}

\begin{region}[pitch=13.2]{274.2}
The \code{blit} method copies an image (more precisely: a \code{Surface}
object) to a specific position on another surface, here the screen. This
lets you display arbitrary images instead of simple shapes.
\end{region}

\begin{region}{322.5}
\emoji{1F4A5} \textbf{Fun fact: where does ``BitBlt'' come from?}
\end{region}

\begin{region}[pitch=13.42]{339.2}
The term \textbf{BitBlt} (bit block transfer) comes from the legendary Xerox
PARC research labs of the 1970s. That is where the graphical user interface
(GUI) as we know it today was developed -- with windows, icons, and mouse
pointers. To display all of this smoothly, a fast method was needed to copy
image data from one memory location to another. That was the birth of the
\textbf{BitBlt function} -- and the word \textbf{blit} has survived to this
day as a short form in many graphics libraries (such as Pygame).
\end{region}

\begin{region}{424.8}
\subheading{Task}
\end{region}

\begin{region}[pitch=13.5]{438.4}
\begin{plaintasks}
\item Test the new feature by running \code{python breakout.py}.
\item As mentioned above, \code{blit()} is a method -- namely of the class
  \code{pygame.Surface}. Both \code{self.img} and \code{Game.screen} are
  objects of this class. Now change the line in the \code{draw} method to:
\end{plaintasks}
\end{region}

\begin{region}[x=74.7,right=0]{494.6}
\tcode{Game.screen.blit(self.img, (self.x, self.y))}
\end{region}

\begin{region}[x=74.7,right=0,pitch=14.27]{515.1}
This does \textbf{exactly the same} as before. Internally, Python more or
less rewrites it into what was there before. In the new form, however, it is
clearer what is meant: \textbf{The blit method is applied to the Game.screen
object to copy the self.img object to the position (self.x, self.y).}
\end{region}

\begin{region}[pitch=13.44]{571.2}
\begin{plaintasks}[start=3]
\item We still have a small cosmetic flaw: the images of the space invaders
  are currently drawn such that \code{(self.x, self.y)} is the top-left
  corner of the image.\\
  Concretely, this means, e.g., that rockets seem to ``fly through'' when
  you hit the image on its right side with a rocket. At the same time, the
  space invaders are already shot down when the rocket actually flies past
  them on the left.\\
  \textbf{Test this!} We don't want to solve problems that we don't know
  exactly.
\end{plaintasks}
\end{region}

\begin{region}[pitch=13.2]{651.9}
\begin{plaintasks}[start=4]
\item To correct this, change the code in the \code{draw} method to:
\end{plaintasks}
\end{region}

\begin{region}[x=74.7,right=0,pitch=13.2]{669.8}
\begin{codeblock}
pos = self.img.get_rect()
pos.center = (self.x, self.y)
Game.screen.blit(self.img, pos)
\end{codeblock}
\end{region}

\begin{region}[x=74.7,right=0,pitch=13.45]{713.8}
The \code{get_rect()} method returns an object of type \code{pygame.Rect}.\\
Such objects represent rectangles (in this case, the border of the image) and
can be centered by overwriting their center attribute.\\
\textbf{Test whether this solves the problem}. If so, then:\\
\code{git commit -a -m "image for space invader"}\\[-0.4bp]
\code{git push}
\end{region}

\nextpage
%% ===================================================================== 4 ==
\begin{region}{82.6}
\heading{Step 2: More images}
\end{region}

\begin{region}[pitch=13.2]{100.6}
Open the file \code{rocket.py} with the class \code{Rocket}, which is
responsible for the rockets. They should now, of course, look like real
rockets!
\end{region}

\begin{region}{140.2}
Task
\end{region}

\begin{region}[pitch=13.4]{153.3}
\begin{tasks}
\item You have seen how it worked for the spaceships. Do the same for the
  rockets and use the image \code{data/rocket0.png} for it.
\item If you transferred the concept directly, the images of the rockets are
  currently aligned at their center.\\
  For rockets, this is actually not quite right -- as with a soldering iron,
  it is the \textbf{tip} that is the dangerous part. Here, however, it is
  harder to test whether a space invader is only hit when it is (currently)
  really caught by the middle of the rocket. So let's briefly throttle the
  game speed a little so that such details can be seen better.\\
  In \code{breakout.py}, change the line:
\end{tasks}
\end{region}

\begin{region}[x=74.7,right=0,pitch=17.7]{291.5}
\tcode{time.sleep(1/30)}\\
to:\\
\tcode{time.sleep(1)}
\end{region}

\begin{region}[x=74.7,right=0,pitch=13.17]{344.6}
Then only one frame per second is drawn -- so we see everything in slow
motion. If you start the game now, you should immediately see that when a
rocket is launched, half of it already sticks out of the paddle (our
``launch pad'').\\
Now, in rocket.py, change the line:
\end{region}

\begin{region}[x=74.7,right=0,pitch=17.75]{401.9}
\tcode{pos.center = (self.x, self.y)}\\
to:\\
\tcode{pos.midtop = (self.x, self.y)}
\end{region}

\begin{region}[x=74.7,right=0,pitch=13.2]{455.0}
If you now restart the game, you will see that the rocket is aligned at its
tip (the dangerous part).
\end{region}

\begin{region}[pitch=13.45]{481.3}
\begin{tasks}[start=3]
\item Since you are currently running the game in slow motion, you may
  notice a problem that was hidden before: the controls are sluggish --
  people say they \textbf{lag} (\textit{lag} = delay).\\
  If you quickly turn the potentiometer back and forth and then let go, the
  paddle still keeps moving back and forth for a while.\\
  This is a typical problem when data is produced and processed at different
  rates (in electronics, this is called \textit{crossing clock domains}).\\
  We will take care of this in the next step.\\
  If everything else works (and you have observed the problem that we will
  solve in the next step), then:\\
  \code{git commit -a -m "rocket image"}\\
  \code{git push}
\end{tasks}
\end{region}

\nextpage
%% ===================================================================== 5 ==
\begin{region}{69.4}
\heading{Step 3: Processing data from the controller faster}
\end{region}

\begin{region}[pitch=13.2]{87.4}
There you will find the code that reads and processes data from the serial
interface (via our module \code{controller}).\\
It should look roughly like this:
\end{region}

\begin{region}[pitch=12]{140.5}
\begin{codeblock}[escapechar=`]
ser_data = controller.readdata()
`\Kw{if}` ser_data:
    `\Kw{if}` ser_data[`\Nu{0}`] == `\Nu{"A"}`:
        rockets.append(Rocket(Paddle.x + Paddle.w // `\Nu{2}`, Paddle.y, -`\Nu{10}`))
    `\Kw{elif}` ser_data[`\Nu{0}`] == `\Nu{"X"}`:
        Paddle.set(ser_data[`\Nu{1}`] * `\Nu{1.5}`)
\end{codeblock}
\end{region}

\begin{region}[pitch=13.8]{226.3}
\textbf{The problem:} This code is executed \textbf{only once per frame},
but the serial interface delivers data \textbf{much faster}.\\
As a result, data ``piles up'', as you saw in the last step.
\end{region}

\begin{region}{280.3}
Fortunately, in our case the problem can be solved very easily by changing
the code to:
\end{region}

\begin{region}[pitch=12]{300.3}
\begin{codeblock}[escapechar=`]
`\Kw{while}` `\Cy{True}`:
    ser_data = controller.readdata()
    `\Kw{if}` ser_data is `\Cy{None}`:
        `\Kw{break}`
    `\Kw{if}` ser_data[`\Nu{0}`] == `\Nu{"A"}`:
        rockets.append(Rocket(Paddle.x + Paddle.w // `\Nu{2}`, Paddle.y, -`\Nu{10}`))
    `\Kw{elif}` ser_data[`\Nu{0}`] == `\Nu{"X"}`:
        Paddle.set(ser_data[`\Nu{1}`] * `\Nu{1.5}`)
\end{codeblock}
\end{region}

\begin{region}[pitch=13.15]{409.2}
This way, the data is processed practically immediately: the loop reads data
until no new data is available and processes it right away -- e.g., by
creating new rockets or setting the paddle position.
\end{region}

\begin{region}{462.8}
\subheading{Task:}
\end{region}

\begin{region}[pitch=13.73]{490.2}
\begin{tasks}
\item Check that with this change the game is \textbf{still very sluggish}
  -- but only because we still draw \textbf{only one frame per second}.\\
  The delays caused by the data processing, however, are gone.\\
  If you now quickly turn the potentiometer and then stop, the paddle is
  drawn \textbf{immediately at the new position} in the next frame and then
  stays there.
\end{tasks}
\end{region}

\begin{region}[pitch=13.3]{558.4}
\begin{tasks}[start=2]
\item Of course, we now want to play with at least 30 frames per second
  again.\\
  But this time we do it properly and no longer use \code{time.sleep()},
  but the Clock that Pygame provides for this purpose.\\
  It is used according to this pattern:
\end{tasks}
\end{region}

\begin{region}[x=74.7,right=0,pitch=13.2]{624.9}
\begin{codeblock}[escapechar=`]
clock = pygame.time.Clock()
`\Kw{while}` running:
    `\Bl{\# here comes the code to handle events}`
    `\Bl{\# and to draw a frame}`

    `\Bl{\# and at the end of the loop:}`
    clock.tick(`\Nu{30}`)
\end{codeblock}
\end{region}

\begin{region}[x=74.7,right=0,pitch=13.1]{743.4}
Apply this pattern and check that the game now runs smoothly again --
possibly even more smoothly than before.
\end{region}

\nextpage
%% ===================================================================== 6 ==
\begin{region}[x=74.7,right=0]{82.6}
\subheading{What is behind this?}
\end{region}

\begin{region}[x=74.7,right=0,pitch=13.3]{97.4}
\begin{dotlist}
\item Before the main loop, a Clock object is created. This object
  internally stores the current time.
\item When you call the method tick(30), it checks:
  \begin{dotlist}
  \item \textbf{If less than 1/30 of a second has passed since the stored
    time}, it blocks briefly until the time span is over, and then stores
    the new current time.
  \item \textbf{If enough time has already passed}, it simply stores the
    current time and continues immediately.
  \end{dotlist}
\end{dotlist}
\end{region}

\begin{region}[x=74.7,right=0,pitch=13.75]{204.4}
\subheading{What is the advantage?}
Drawing does not always take exactly the same amount of time!\\
How long it takes depends, e.g., on how many space invaders or balls are
currently in the game, or on whether an object has just been hit.\\
As a rule, however, drawing takes considerably \textbf{less} than 1/30 of a
second.\\
This mechanism ensures that a frame is always redrawn \textbf{after the same
time span} -- regardless of how much is going on.
\end{region}

\begin{region}[pitch=13.45]{300.1}
\begin{tasks}[start=3]
\item If everything works, into the repository with it:\\
  \code{git commit -a -m "fixed lag and used clock"}\\
  \code{git push}
\end{tasks}
\end{region}

\begin{region}{383.8}
\heading{Step 4: Explosions \emoji{1F4A5}}
\end{region}

\begin{region}[pitch=13.3]{404.0}
After the boring data processing, we need some action -- the space invaders
should explode when they are hit!\\
As luck would have it, \code{data/explosion.png} already provides a suitable
image for this.
\end{region}

\begin{region}[pitch=13.5]{457.3}
The first approach is to check in the \code{draw()} method of the Ship class
whether the attribute \code{alive} has the value \code{True} or
\code{False}.\\
If it is \code{True}, the ``normal'' image is displayed -- otherwise the one
with the explosion.\\
Open the file \code{ship.py} and add the following line in the constructor
(i.e., the \code{__init__} method):
\end{region}

\begin{region}{515.3}
\begin{codeblock}[escapechar=`]
self.die = pygame.image.load(`\Nu{"data/explosion.png"}`)
\end{codeblock}
\end{region}

\begin{region}{544.9}
Then change the \code{draw} method as follows:
\end{region}

\begin{region}[pitch=12]{571.5}
\begin{codeblock}[escapechar=`]
`\Kw{def}` `\Cy{draw}\Bl{(self):}`
    `\Kw{if}` self.alive:
        img = self.img
    `\Kw{else}`:
        img = self.die
    pos = img.get_rect()
    pos.center = (self.x, self.y)
    Game.screen.blit(img, pos)
\end{codeblock}
\end{region}

\begin{region}{693.4}
\subheading{Task}
\end{region}

\begin{region}[pitch=14.1]{718.7}
\begin{tasks}
\item Make the change and test whether it works. You should notice that the
  program runs \textbf{without errors}, but \textbf{no explosions are
  displayed}.
\end{tasks}
\end{region}

\begin{region}[pitch=13.55]{747.0}
\begin{tasks}[start=2]
\item \textbf{Why are no explosions displayed?}\\
  The problem is: a spaceship is no longer drawn at all when its
  \code{alive} attribute is \code{False}.
\end{tasks}
\end{region}

\nextpage
%% ===================================================================== 7 ==
\begin{region}[x=74.7,right=0,pitch=13.2]{68.8}
Have a look at the file \code{breakout.py} -- in the main loop, there is a
line like this one:
\end{region}

\begin{region}[x=74.7,right=0]{95.4}
\begin{codeblock}[escapechar=`]
ships = [item `\Kw{for}` item `\Kw{in}` ships `\Kw{if}` item.alive]
\end{codeblock}
\end{region}

\begin{region}[x=74.7,right=0,pitch=13.75]{122.1}
Here, \code{ships} is the list of spaceships that are drawn. When a ship has
been hit, it is removed from the list \textbf{immediately} -- and therefore
no longer appears for a last ``explosion frame''.
\end{region}

\begin{region}[x=74.7,right=0]{176.9}
\subheading{This is how we solve the problem:}
\end{region}

\begin{region}[x=74.7,right=0,pitch=14.2]{204.4}
\begin{boldsteps}
\item \textbf{Marking hits differently}\\
  In \code{breakout.py}, find the place where a hit has so far been handled with:
\end{boldsteps}
\end{region}

\begin{region}[x=92.7,right=0,pitch=17.75]{236.6}
\tcode{s.alive = }{\ttfamily\Cy{False}}\\
Change it to:\\
\tcode{s.hit = }{\ttfamily\Cy{True}}
\end{region}

\begin{region}[x=92.7,right=0,pitch=13.2]{290.7}
\textbf{Tip}: Test right away whether the program still runs without errors
-- that way you catch typos early. But don't be surprised: for now, the space
invaders are invulnerable -- we are not done yet!
\end{region}

\begin{region}[x=74.7,right=0,pitch=13.5]{331.3}
\begin{boldsteps}[start=2]
\item \textbf{Initializing and using the \codearg{\mdseries hit}
  attribute}\\
  In \code{ship.py}, initialize the new attribute in the constructor:
\end{boldsteps}
\end{region}

\begin{region}[x=92.7,right=0,pitch=17.9]{362.7}
\tcode{self.hit = }{\ttfamily\Cy{False}}\\
In the \code{draw} method, change the code as follows:
\end{region}

\begin{region}[x=96.7,right=0,pitch=12]{398.2}
\begin{codeblock}[escapechar=`]
`\Kw{def}` `\Cy{draw}\Bl{(self):}`
    `\Kw{if}` self.hit:
        img = self.die
        self.alive = `\Cy{False}`
    `\Kw{else}`:
        img = self.img
    pos = img.get_rect()
    pos.center = (self.x, self.y)
    Game.screen.blit(img, pos)
\end{codeblock}
\end{region}

\begin{region}[x=92.7,right=0,pitch=13.7]{517.8}
Test whether explosions are displayed now!\\
A space invader that has been hit should remain visible as an explosion for
\textbf{exactly one frame} before it disappears from the list (and from
life).
\end{region}

\begin{region}[pitch=15.5]{562.7}
\begin{plaintasks}[start=3]
\item Now that a solution has been spoon-fed to you, you may solve something
  on your own again! \emoji{1F604}\\[-1.2bp]
  When a rocket hits a space invader, \textbf{the rocket should explode,
  too}.\\[-1.2bp]
  When a ball hits the space invader, the ball should \textbf{not} explode.
\end{plaintasks}
\end{region}

\begin{region}[x=74.7,right=0]{621.1}
\textbf{A possible approach (analogous to the previous subtask)}:
\end{region}

\begin{region}[x=74.7,right=0,pitch=14.1]{635.4}
\begin{steps}
\item In \code{breakout.py}, the attribute hit is set not only for the ship
  that was hit, \textbf{but also for the object that caused the hit}.
\item In the \code{Ball} class, the \code{hit} attribute is ignored.\\
  In the \code{Rocket} class, it is \textbf{handled the same way as in the
  Ship class}.
\end{steps}
\end{region}

\begin{region}[pitch=13.3]{690.8}
\begin{plaintasks}[start=4]
\item If everything works, into the repository with it:\\
  \code{git commit -a -m "explosions"}\\
  \code{git push}
\end{plaintasks}
\end{region}

\nextpage
%% ===================================================================== 8 ==
\begin{region}{69.4}
\heading{Step 5: Cleaning up the code}
\end{region}

\begin{region}[pitch=13.5]{87.4}
Have a look at the file \code{wall.py}. Currently, it is rather ugly that
things like the width and height of the bricks and also the maximum number
of bricks per row and column appear \textbf{hard-coded} as constants.\\
You probably feel the same way I did: I had to look at it for a while to
understand what exactly was specified. So I have decoded it for you:
\end{region}

\begin{region}[pitch=14.2]{154.4}
\begin{dashlist}
\item A brick is \textbf{80 pixels wide} and \textbf{20 pixels high}.
\item There are at most \textbf{10 brick columns} and \textbf{30 brick
  rows}.
\end{dashlist}
\end{region}

\begin{region}[pitch=13.7]{183.2}
Also have a look at the file \code{game.py} in the other terminal.\\
There, the size of the drawing area is also hard-coded in the \code{init}
method:\\
\textbf{800 pixels wide}, \textbf{600 pixels high}.\\
Fortunately, this fits together as it is currently specified:
\end{region}

\begin{region}[pitch=14.7]{237.5}
\begin{dashlist}
\item 80 pixels \texttimes\ 10 columns $\rightarrow$ 800 pixels wide.
\item 20 pixels \texttimes\ 30 rows $\rightarrow$ 600 pixels high.
\end{dashlist}
\end{region}

\begin{region}[pitch=13.33]{272.0}
\subheading{Great, so everything is fine, right?}
Well, yes. But what if you want to change the size of the bricks? Or want
more bricks?\\
Or what if you -- like me -- still want to understand the code in a few
weeks?\\
Let's organize it better! Change \code{game.py} as follows:
\end{region}

\begin{region}{329.9}
\begin{codeblock}[escapechar=`]
`\Mg{import}` `\Bl{pygame}`
\end{codeblock}
\end{region}

\begin{region}[pitch=12]{355.1}
\begin{codeblock}[escapechar=`]
`\Kw{class}` `\Cy{Game}`:
    width, height = `\Cy{None}`, `\Cy{None}`
    screen = `\Cy{None}`
    brick_rows, brick_cols = `\Nu{30}`, `\Nu{10}`
    brick_width, brick_height = `\Nu{80}`, `\Nu{20}`

    `\Kw{def}` `\Cy{init}`():
        Game.width = Game.brick_cols * Game.brick_width
        Game.height = Game.brick_rows * Game.brick_height
        pygame.init()
        Game.screen = pygame.display.set_mode((Game.width, Game.height))

    `\Kw{def}` `\Cy{quit}`():
        pygame.quit()
\end{codeblock}
\end{region}

\begin{region}[pitch=13.2]{527.2}
Now it is specified in one central place how large the bricks are and how
many rows and columns there are at most.
\end{region}

\begin{region}{558.9}
\subheading{Task}
\end{region}

\begin{region}[pitch=13.5]{576.8}
\begin{tasks}
\item Now it's your turn to make use of the new \code{Game} attributes in
  \code{wall.py}:
  \begin{dotlist}
  \item Replace every occurrence of \code{30} with \code{Game.brick_rows}.
  \item Replace every occurrence of \code{10} with \code{Game.brick_cols}.
  \item Replace every occurrence of \code{80} with \code{Game.brick_width}.
  \item Replace every occurrence of \code{20} with
    \code{Game.brick_height}.
  \end{dotlist}
  Then test whether everything still works!
\end{tasks}
\end{region}

\begin{region}[pitch=13.4]{657.5}
\begin{tasks}[start=2]
\item Also check whether you can change the new constants in \code{game.py}
  and everything still works. For example, change \code{Game.brick_cols} to
  \code{15}. Test whether the game still works then.
\end{tasks}
\end{region}

\begin{region}[pitch=17.5]{697.5}
\begin{tasks}[start=3]
\item Afterwards, undo the change -- a drawing area of 800 \texttimes\ 600
  pixels is not bad if you want to play the game on a Raspberry Pi later
  on. \emoji{1F604}
\end{tasks}
\end{region}

\begin{region}[pitch=13.4]{731.7}
\begin{tasks}[start=4]
\item If everything works:\\
  \code{git commit -a -m "brick constants"}\\[-0.2bp]
  \code{git push}
\end{tasks}
\end{region}

\nextpage
%% ===================================================================== 9 ==
\begin{region}{82.6}
\heading{Step 6: Even more code cleanup}
\end{region}

\begin{region}[pitch=13.2]{100.6}
Have a look at the file \code{breakout.py}.\\
The main loop is slowly getting quite large, because at the moment
everything happens in it at once:
\begin{dashlist}
\item We poll the serial interface for new data.
\item We poll Pygame for new key events.
\item We draw and move balls, spaceships, and rockets.
\item We check whether there are collisions.
\end{dashlist}
\end{region}

\begin{region}[pitch=13.3]{193.8}
Now at least the \textbf{drawing and moving} of balls, spaceships, and
rockets should be organized better.
\end{region}

\begin{region}[pitch=13.23]{234.4}
\subheading{Why?}
\begin{boldsteps}
\item The code is almost identical.
\item When the spaceships also drop bombs later, another list of objects
  will be added that have to be moved and drawn.
\end{boldsteps}
\end{region}

\begin{region}{301.4}
So: \textbf{Now is the right time to do this properly!}
\end{region}

\begin{region}[pitch=13.37]{327.8}
And we'll do it completely right straight away:\\
The lists of the objects that are moved and drawn -- i.e., \code{balls},
\code{rockets}, and \code{ships} -- are moved into the Game class.\\
There will also be a method \code{draw()} there that manages all of these
lists.
\end{region}

\begin{region}{394.6}
Change \code{game.py} as follows:
\end{region}

\begin{region}{421.3}
\begin{codeblock}[escapechar=`]
`\Mg{import}` `\Bl{pygame}`
\end{codeblock}
\end{region}

\begin{region}[pitch=12]{446.5}
\begin{codeblock}[escapechar=`]
`\Kw{class}` `\Cy{Game}`:
    width, height = `\Cy{None}`, `\Cy{None}`
    screen = `\Cy{None}`
    brick_rows = `\Nu{30}`
    brick_cols = `\Nu{15}`
    brick_width = `\Nu{80}`
    brick_height = `\Nu{20}`
    balls = []
    rockets = []
    ships = []

    `\Kw{def}` `\Cy{init}`():
        Game.width = Game.brick_cols * Game.brick_width
        Game.height = Game.brick_rows * Game.brick_height
        pygame.init()
        Game.screen = pygame.display.set_mode((Game.width, Game.height))

    `\Kw{def}` `\Cy{draw}`():
        `\Kw{for}` item `\Kw{in}` Game.balls + Game.rockets + Game.ships:
            item.draw()
            item.move()
        Game.balls = [item `\Kw{for}` item `\Kw{in}` Game.balls `\Kw{if}` item.alive]
        Game.rockets = [item `\Kw{for}` item `\Kw{in}` Game.rockets `\Kw{if}` item.alive]
        Game.ships = [item `\Kw{for}` item `\Kw{in}` Game.ships `\Kw{if}` item.alive]

    `\Kw{def}` `\Cy{quit}`():
        pygame.quit()
\end{codeblock}
\end{region}

\nextpage
%% ==================================================================== 10 ==
\begin{region}{95.8}
\subheading{Tasks}
\end{region}

\begin{region}[pitch=13.55]{122.5}
\begin{tasks}
\item Adapt \code{breakout.py} so that the change takes effect:
  \begin{dotlist}
  \item Remove the code that has so far moved and drawn the balls,
    spaceships, and rockets.\\
    Instead, \code{Game.draw()} should now be called once.
  \item All places where the lists \code{balls}, \code{rockets}, or
    \code{ships} are still accessed directly must now be adapted as
    well.\\
    There, write, e.g., \code{Game.balls} instead of \code{balls}, etc.
  \item Then test whether everything works again!
  \end{dotlist}
\end{tasks}
\end{region}

\begin{region}[pitch=13.8]{217.0}
\begin{tasks}[start=2]
\item To see that this cleanup is useful, too, make another small change in
  \code{wall.py}: whenever a brick is hit by a ball, a new space invader
  should enter the game. \textbf{Sounds complicated?} Don't worry, this
  \textbf{takes just two lines}:
\end{tasks}
\end{region}

\begin{region}[x=74.7,right=0,pitch=13.3]{258.2}
\begin{boldsteps}
\item At the beginning of the file, import the Ship class with
  \code{from ship import Ship}
\item When a brick has been hit, append a new Ship object to the list
  Game.ships. To do this, find the line with \code{Wall.brick[i2, j2] = 0}
  (which vaporizes a brick that has been hit) and insert the line
\end{boldsteps}
\end{region}

\begin{region}[x=92.7,right=0,pitch=17.7]{315.9}
\tcode{Game.ships.append(Ship(0))}\\
below it.
\end{region}

\begin{region}[x=74.7,right=0,pitch=13.3]{352.1}
\textbf{Test whether this works:} every destroyed brick should now bring a
new spaceship into the game.
\end{region}

\begin{region}[x=74.7,right=0,pitch=13.4]{383.9}
\subheading{If everything works:}
\code{git commit -a -m "moving stuff is in Game"}\\
\code{git push}
\end{region}

\nextpage
%% ==================================================================== 11 ==
\begin{region}{69.4}
\heading{Step 7: A completely different construction site}
\end{region}

\begin{region}[pitch=13.38]{88.1}
Test what happens when you \textbf{unplug the microcontroller while the game
is running}.\\
The program then crashes.\\
According to the error message, the problem occurs in \code{controller.py},
more precisely in the method \code{readline()}.\\
It currently has this structure:
\end{region}

\begin{region}[pitch=12]{168.2}
\begin{codeblock}[escapechar=`]
`\Kw{def}` `\Cy{readline}`():
    `\Kw{if}` Controller.ser `\Kw{is}` `\Cy{None}`:
        `\Kw{return}` `\Cy{None}`

    `\Kw{while}` Controller.ser.in_waiting > `\Nu{0}`:
        `\Cm{\# code to read the serial interface}`

    `\Kw{return}` `\Cy{None}`
\end{codeblock}
\end{region}

\begin{region}[pitch=13.53]{292.1}
\subheading{What happens here?}
If the microcontroller is removed while the game is running,
\code{Controller.ser} is not \code{None}, but:\\
as soon as you access the object (e.g., \code{Controller.ser.in_waiting}), it
crashes.
\end{region}

\begin{region}[pitch=13.47]{360.0}
\subheading{The solution:}
We put the while loop in which the object is used into a \code{try} block --
and in the corresponding \code{except} block we catch the error by setting
\code{Controller.ser} back to \code{None}.\\
So the structure is adapted like this:
\end{region}

\begin{region}{440.5}
\begin{codeblock}[escapechar=`]
`\Kw{def}` `\Cy{readline}\Bl{():}`
\end{codeblock}
\end{region}

\begin{region}[pitch=12]{453.7}
\begin{codeblock}[escapechar=`]
    `\Kw{if}` Controller.ser `\Kw{is}` `\Cy{None}`:
        `\Kw{return}` `\Cy{None}`

    `\Kw{try}`:
        `\Kw{while}` Controller.ser.in_waiting > `\Nu{0}`:
            `\Cm{\# code as before}`

    `\Kw{except}`:
        `\Cy{print}`(`\Nu{"disconnected"}`)
        Controller.ser = `\Cy{None}`
    `\Kw{return}` `\Cy{None}`
\end{codeblock}
\end{region}

\begin{region}{612.7}
\subheading{Task}
\end{region}

\begin{region}[pitch=13.2]{639.2}
Make this adaptation and test whether you can now remove the microcontroller
while the game is running without the program crashing.
\end{region}

\begin{region}[pitch=13.55]{679.6}
\textbf{Note}: By the way, if you have not put \code{ser} into a class
Controller (i.e., your code contains \code{ser.in_waiting} instead of
\code{Controller.ser.in_waiting}),\\
then the \code{except} block should of course simply contain
\code{ser = None}.
\end{region}

\begin{region}[pitch=13.3]{734.1}
\textbf{Remember}: When you reconnect the microcontroller, you have to run
usb\_check on the lab computers so that it is recognized again.
\end{region}

\nextpage
%% ==================================================================== 12 ==
\begin{region}{69.4}
\heading{Step 8: A few bombs to finish}
\end{region}

\begin{region}[pitch=14]{87.2}
The poor space invaders cannot defend themselves at all. In this step, we
allow them to \textbf{drop bombs}. It won't really help them today, though
-- the bombs won't destroy our paddle (yet).\\
The main point here is to test \textbf{how easily a new feature can be built
into the reorganized code}. And to show you \textbf{how to implement
something like this quick and dirty, step by step}.
\end{region}

\begin{region}{175.9}
\subheading{Task}
\end{region}

\begin{region}{193.9}
\begin{plaintasks}
\item At the very top of \code{ship.py}, import:
\end{plaintasks}
\end{region}

\begin{region}[x=74.7,right=0,pitch=13.2]{211.8}
\begin{codeblock}
import random
from rocket import Rocket
\end{codeblock}
\end{region}

\begin{region}[x=74.7,right=0,pitch=13.3]{242.6}
We need the module \code{random} for random numbers, and with \code{Rocket}
we get our class for rockets.\\
In the \code{move()} method of the Ship class, add these two lines:
\end{region}

\begin{region}[x=74.7,right=0,pitch=13.2]{287.1}
\begin{codeblock}
if random.randint(1, 100) == 1:
    Game.rockets.append(Rocket(self.x, self.y, 10))
\end{codeblock}
\end{region}

\begin{region}[x=74.7,right=0,pitch=13.4]{318.6}
\textbf{What happens here?} In every frame, each spaceship has a 1\% chance
of dropping a ``bomb''. Okay, strictly speaking, it is not a bomb but a
rocket -- but those are details that can be changed later. \emoji{1F604}
\end{region}

\begin{region}[x=74.7,right=0,pitch=13.2]{369.6}
Test what happens! You should observe: the spaceships really do drop rockets
downward (which, for now, still point in the wrong direction). But in this
``attack'', they currently destroy themselves.
\end{region}

\begin{region}[pitch=13.57]{410.1}
\begin{plaintasks}[start=2]
\item \textbf{Why do the spaceships destroy themselves?} This happens
  because currently, for every object in Game.rockets, it is checked whether
  the distance to a spaceship is small enough for a contact. \textbf{The
  solution}: we create a new list Game.bombs in Game, so that bombs are
  managed separately from rockets.
\end{plaintasks}
\end{region}

\begin{region}[x=74.7,right=0,pitch=13.5]{464.3}
\begin{steps}
\item Add a new, empty list \code{bombs} to \code{Game}, analogous to the
  existing lists \code{balls}, \code{rockets}, and \code{ships}.
\item Change the \code{draw} method of \code{Game} so that
  \begin{dotlist}
  \item all objects in \code{Game.balls}, \code{Game.rockets},
    \code{Game.ships}, and now also in \code{Game.bombs} are moved and
    drawn.
  \item objects that are no longer alive are removed from
    \code{Game.bombs}.
  \end{dotlist}
\item Change the bomb dropping in \code{ship.py} so that the rockets are
  appended to \code{Game.bombs}.
\end{steps}
\end{region}

\begin{region}[x=74.7,right=0]{571.7}
Then test the game! The space invaders should now survive their bomb
dropping.
\end{region}

\begin{region}[pitch=13.2]{584.9}
\begin{plaintasks}[start=3]
\item Of course, it is still unattractive that rockets are currently simply
  dropped instead of bombs. No problem -- by now you know how to solve
  something like this analogously!
\end{plaintasks}
\end{region}

\begin{region}[x=74.7,right=0,pitch=13.5]{615.9}
Use \code{cp rocket.py bomb.py} to create a copy of \code{rocket.py}.\\
In \code{bomb.py}, you then make the following changes:
\begin{steps}
\item Name the class \code{Bomb}.
\item Use \code{data/bomb.png} as the image.
\item In the \code{move} method, make sure that \code{self.alive} is set to
  \code{False}\\ if \code{self.y > Game.height}.
\end{steps}
And in \code{ship.py}, you now import the class \code{Bomb} instead of
\code{Rocket}, and you also append a \code{Bomb} object instead of a
\code{Rocket} object.\\
Test whether bombs really fly now! If all is well, your code may go into
the repository:
\end{region}

\begin{region}[x=74.7,right=0,pitch=13.35]{741.5}
\code{git add bomb.py}\\
\code{git commit -a -m "bomb blast"}\\
\code{git push}
\end{region}

\end{document}
